\section{Estat de l'art}

\subsection{RCA tradicional i limitacions en entorns data-rich}

Les metodologies clàssiques de RCA (5 Whys, Ishikawa, FMEA) han estat àmpliament utilitzades en la indústria farmacèutica per estructurar investigacions de desviació. Tanmateix, diversos treballs evidencien limitacions quan el volum i heterogeneïtat de dades augmenten: dependència elevada del coneixement tàcit, baixa repetibilitat entre equips i dificultat per capturar interaccions multivariants complexes.

En entorns GxP, l'RCA requereix evidència auditable i justificació de decisions. Per això, l'objectiu no és automatitzar la decisió final, sinó reforçar la capacitat d'anàlisi de l'equip investigador amb eines quantitatives que prioritzin hipòtesis de manera transparent.

\subsection{Machine Learning i explicabilitat per a qualitat de processos}

En els darrers anys, l'aplicació de ML en qualitat farmacèutica ha crescut en casos com detecció d'anomalies, manteniment predictiu i predicció de lots fora d'especificació, utilitzant models com \emph{random forest}, \emph{gradient boosting} \cite{xgboost} o xarxes neuronals. Tanmateix, en escenaris RCA regulats, una bona mètrica predictiva és condició necessària però no suficient: cal demostrar robustesa, estabilitat temporal i, especialment, interpretabilitat operativa \cite{doshi2017rigorous,rudin2021fundamental}.

Les tècniques XAI (SHAP \cite{shap_arxiv1705}, LIME \cite{lime_arxiv1602}) han permès avançar en la transformació de prediccions en hipòtesis investigables, actuant com a mecanisme de \textbf{priorització d'hipòtesis} per reduir l'espai de recerca. Diversos treballs destaquen la diferència entre models orientats al rendiment i models interpretable-by-design que, amb compromís moderat en exactitud, faciliten la justificació de resultats i adopció en processos de qualitat.

\subsection{Explainable Boosting Machines: estat actual}

Els \emph{Explainable Boosting Machines} (EBM) representen una extensió moderna dels \emph{Generalized Additive Models} (GAMs) \cite{hastie1990_gam} que combina interpretabilitat nativa amb capacitat predictiva competitiva. Els treballs de Lou et al. \cite{lou2013_ga2m} i Caruana et al. \cite{caruana2015_intelligible} van establir les bases dels models GA2M (GAMs amb interaccions de parell), demostrant que és possible capturar interaccions controlades mantenint una estructura additiva llegible.

La literatura recent mostra que els EBM assoleixen rendiment comparable a \emph{Random Forest} o \emph{Boosted Trees}, amb l'avantatge d'un funcionament intern directament explicable \cite{ebm_nori,piml_ebm,ebm_2025}. A diferència dels models \emph{black box}, l'EBM pertany a la família \emph{glass box} i permet descriure explícitament la contribució de cada variable i interacció a la predicció, facilitant la comunicació amb experts de domini i la seva adopció en contextos on la transparència és crítica.

\subsection{Industrialització i validació en entorns GxP}

La literatura de regulació digital i validació de sistemes computaritzats (GAMP 5 \cite{gamp5}, ALCOA+ \cite{alcoa}) estableix la necessitat de definir l'ús previst, control de canvis, traçabilitat de dades/models i monitoratge continu del rendiment per garantir que una solució ML sigui defensable en auditoria. Diversos treballs destaquen també el risc de \emph{drift} \cite{lu2020drift} i la necessitat de governança del cicle de vida del model.

Com a síntesi, l'estat de l'art mostra que existeixen peces madures (RCA estructurat, modelització predictiva, XAI, qualitat de dades), però manca sovint una metodologia integral que les uneixi de forma operativa i compatible amb entorns regulats. Aquesta és precisament la bretxa que aborda el present treball: utilitzar un model interpretable-by-design com l'EBM per prioritzar hipòtesis de causa arrel amb criteris tècnics, traçables i alineats amb requisits GxP.